\documentclass[a4paper,12pt,twoside]{memoir}

% Castellano
\usepackage[spanish,es-tabla]{babel}
\selectlanguage{spanish}
\usepackage[utf8]{inputenc}
\usepackage[T1]{fontenc}
\usepackage{lmodern} % scalable font
\usepackage{microtype}
\usepackage{placeins}

\RequirePackage{booktabs}
\RequirePackage[table]{xcolor}
\RequirePackage{xtab}
\RequirePackage{multirow}

% Links
\PassOptionsToPackage{hyphens}{url}\usepackage[colorlinks]{hyperref}
\hypersetup{
	allcolors = {red}
}

% Ecuaciones
\usepackage{amsmath}

% Rutas de fichero / paquete
\newcommand{\ruta}[1]{{\sffamily #1}}

% Párrafos
\nonzeroparskip

% Huérfanas y viudas
\widowpenalty100000
\clubpenalty100000

% Evitar solapes en el header
\nouppercaseheads

% Imagenes
\usepackage{graphicx}
\newcommand{\imagen}[2]{
	\begin{figure}[!h]
		\centering
		\includegraphics[width=0.9\textwidth]{#1}
		\caption{#2}\label{fig:#1}
	\end{figure}
	\FloatBarrier
}

\newcommand{\imagenflotante}[2]{
	\begin{figure}%[!h]
		\centering
		\includegraphics[width=0.9\textwidth]{#1}
		\caption{#2}\label{fig:#1}
	\end{figure}
}



% El comando \figura nos permite insertar figuras comodamente, y utilizando
% siempre el mismo formato. Los parametros son:
% 1 -> Porcentaje del ancho de página que ocupará la figura (de 0 a 1)
% 2 --> Fichero de la imagen
% 3 --> Texto a pie de imagen
% 4 --> Etiqueta (label) para referencias
% 5 --> Opciones que queramos pasarle al \includegraphics
% 6 --> Opciones de posicionamiento a pasarle a \begin{figure}
\newcommand{\figuraConPosicion}[6]{%
  \setlength{\anchoFloat}{#1\textwidth}%
  \addtolength{\anchoFloat}{-4\fboxsep}%
  \setlength{\anchoFigura}{\anchoFloat}%
  \begin{figure}[#6]
    \begin{center}%
      \Ovalbox{%
        \begin{minipage}{\anchoFloat}%
          \begin{center}%
            \includegraphics[width=\anchoFigura,#5]{#2}%
            \caption{#3}%
            \label{#4}%
          \end{center}%
        \end{minipage}
      }%
    \end{center}%
  \end{figure}%
}

%
% Comando para incluir imágenes en formato apaisado (sin marco).
\newcommand{\figuraApaisadaSinMarco}[5]{%
  \begin{figure}%
    \begin{center}%
    \includegraphics[angle=90,height=#1\textheight,#5]{#2}%
    \caption{#3}%
    \label{#4}%
    \end{center}%
  \end{figure}%
}
% Para las tablas
\newcommand{\otoprule}{\midrule [\heavyrulewidth]}
%
% Nuevo comando para tablas pequeñas (menos de una página).
\newcommand{\tablaSmall}[5]{%
 \begin{table}
  \begin{center}
   \rowcolors {2}{gray!35}{}
   \begin{tabular}{#2}
    \toprule
    #4
    \otoprule
    #5
    \bottomrule
   \end{tabular}
   \caption{#1}
   \label{tabla:#3}
  \end{center}
 \end{table}
}

%
%Para el float H de tablaSmallSinColores
\usepackage{float}

%
% Nuevo comando para tablas pequeñas (menos de una página).
\newcommand{\tablaSmallSinColores}[5]{%
 \begin{table}[H]
  \begin{center}
   \begin{tabular}{#2}
    \toprule
    #4
    \otoprule
    #5
    \bottomrule
   \end{tabular}
   \caption{#1}
   \label{tabla:#3}
  \end{center}
 \end{table}
}

\newcommand{\tablaApaisadaSmall}[5]{%
\begin{landscape}
  \begin{table}
   \begin{center}
    \rowcolors {2}{gray!35}{}
    \begin{tabular}{#2}
     \toprule
     #4
     \otoprule
     #5
     \bottomrule
    \end{tabular}
    \caption{#1}
    \label{tabla:#3}
   \end{center}
  \end{table}
\end{landscape}
}

%
% Nuevo comando para tablas grandes con cabecera y filas alternas coloreadas en gris.
\newcommand{\tabla}[6]{%
  \begin{center}
    \tablefirsthead{
      \toprule
      #5
      \otoprule
    }
    \tablehead{
      \multicolumn{#3}{l}{\small\sl continúa desde la página anterior}\\
      \toprule
      #5
      \otoprule
    }
    \tabletail{
      \hline
      \multicolumn{#3}{r}{\small\sl continúa en la página siguiente}\\
    }
    \tablelasttail{
      \hline
    }
    \bottomcaption{#1}
    \rowcolors {2}{gray!35}{}
    \begin{xtabular}{#2}
      #6
      \bottomrule
    \end{xtabular}
    \label{tabla:#4}
  \end{center}
}

%
% Nuevo comando para tablas grandes con cabecera.
\newcommand{\tablaSinColores}[6]{%
  \begin{center}
    \tablefirsthead{
      \toprule
      #5
      \otoprule
    }
    \tablehead{
      \multicolumn{#3}{l}{\small\sl continúa desde la página anterior}\\
      \toprule
      #5
      \otoprule
    }
    \tabletail{
      \hline
      \multicolumn{#3}{r}{\small\sl continúa en la página siguiente}\\
    }
    \tablelasttail{
      \hline
    }
    \bottomcaption{#1}
    \begin{xtabular}{#2}
      #6
      \bottomrule
    \end{xtabular}
    \label{tabla:#4}
  \end{center}
}

%
% Nuevo comando para tablas grandes sin cabecera.
\newcommand{\tablaSinCabecera}[5]{%
  \begin{center}
    \tablefirsthead{
      \toprule
    }
    \tablehead{
      \multicolumn{#3}{l}{\small\sl continúa desde la página anterior}\\
      \hline
    }
    \tabletail{
      \hline
      \multicolumn{#3}{r}{\small\sl continúa en la página siguiente}\\
    }
    \tablelasttail{
      \hline
    }
    \bottomcaption{#1}
  \begin{xtabular}{#2}
    #5
   \bottomrule
  \end{xtabular}
  \label{tabla:#4}
  \end{center}
}



\definecolor{cgoLight}{HTML}{EEEEEE}
\definecolor{cgoExtralight}{HTML}{FFFFFF}

%
% Nuevo comando para tablas grandes sin cabecera.
\newcommand{\tablaSinCabeceraConBandas}[5]{%
  \begin{center}
    \tablefirsthead{
      \toprule
    }
    \tablehead{
      \multicolumn{#3}{l}{\small\sl continúa desde la página anterior}\\
      \hline
    }
    \tabletail{
      \hline
      \multicolumn{#3}{r}{\small\sl continúa en la página siguiente}\\
    }
    \tablelasttail{
      \hline
    }
    \bottomcaption{#1}
    \rowcolors[]{1}{cgoExtralight}{cgoLight}

  \begin{xtabular}{#2}
    #5
   \bottomrule
  \end{xtabular}
  \label{tabla:#4}
  \end{center}
}




\graphicspath{ {./img/} }

% Capítulos
\chapterstyle{bianchi}
\newcommand{\capitulo}[2]{
	\setcounter{chapter}{#1}
	\setcounter{section}{0}
	\setcounter{figure}{0}
	\setcounter{table}{0}
	\chapter*{#2}
	\addcontentsline{toc}{chapter}{#2}
	\markboth{#2}{#2}
}

% Apéndices
\renewcommand{\appendixname}{Apéndice}
\renewcommand*\cftappendixname{\appendixname}

\newcommand{\apendice}[1]{
	%\renewcommand{\thechapter}{A}
	\chapter{#1}
}

\renewcommand*\cftappendixname{\appendixname\ }

% Formato de portada
\makeatletter
\usepackage{xcolor}
\newcommand{\tutor}[1]{\def\@tutor{#1}}
\newcommand{\course}[1]{\def\@course{#1}}
\definecolor{cpardoBox}{HTML}{E6E6FF}
\def\maketitle{
  \null
  \thispagestyle{empty}
  % Cabecera ----------------
\noindent\includegraphics[width=\textwidth]{cabecera}\vspace{1cm}%
  \vfill
  % Título proyecto y escudo informática ----------------
  \colorbox{cpardoBox}{%
    \begin{minipage}{.8\textwidth}
      \vspace{.5cm}\Large
      \begin{center}
      \textbf{TFM del Máster Universitario en Ingeniería Informática}\vspace{.6cm}\\
      \textbf{\LARGE\@title{}}
      \end{center}
      \vspace{.2cm}
    \end{minipage}

  }%
  \hfill\begin{minipage}{.20\textwidth}
    \includegraphics[width=\textwidth]{escudoInfor}
  \end{minipage}
  \vfill
  % Datos de alumno, curso y tutores ------------------
  \begin{center}%
  {%
    \noindent\LARGE
    Presentado por \@author{}\\ 
    en Universidad de Burgos --- \@date{}\\
    Tutor: \@tutor{}\\
  }%
  \end{center}%
  \null
  \cleardoublepage
  }
\makeatother

\newcommand{\nombre}{Francisco Martín Vargas} %%% cambio de comando
\newcommand{\nombretutor}{Daniel Urda Muñoz} 
\newcommand{\titulo}{Modelo predictivo de resultados de carreras de Fórmula 1} 

% Datos de portada
\title{\titulo \\Documentación Técnica}
\author{\nombre}
\tutor{\nombretutor}
\date{\today}

\begin{document}

\maketitle



\cleardoublepage



%%%%%%%%%%%%%%%%%%%%%%%%%%%%%%%%%%%%%%%%%%%%%%%%%%%%%%%%%%%%%%%%%%%%%%%%%%%%%%%%%%%%%%%%



\frontmatter


\clearpage

% Indices
\tableofcontents

\clearpage

\listoffigures

\clearpage

\listoftables

\clearpage

\mainmatter

\appendix

\apendice{Plan de Proyecto Software}

\section{Introducción}

En el presente anexo se describirá la planificación llevada a cabo en el proyecto. La planificación es una etapa crucial en la cual se estiman los costos del proyecto, tanto en términos de trabajo, tiempo y recursos financieros. Analizar la planificación resulta sumamente beneficioso, puesto que permite determinar la viabilidad del proyecto, el tiempo necesario para su culminación, los recursos requeridos y su cronograma de empleo, entre otros aspectos relevantes. Para ello, se ha divido la planificación en dos fases fundamentales: la planificación temporal y el estudio de viabilidad.

\section{Planificación temporal}

Se va a seguir una planificación por \textit{sprints}, la cual se explica a continuación.

\subsection{Sprint 0: Planificación y Objetivos del Proyecto (2 semana)}
La planificación se llevo a cabo en 3 fases.
\begin{itemize}
    \item
    Planificar el desarrollo en sprints del proyecto.
    La planificación del desarrollo del proyecto comienza con la definición de los objetivos del proyecto y los resultados esperados. Además en este punto vamos a dividir el trabajo en tareas mas pequeñas.
    \item
    Establecer los objetivos principales del proyecto.
    Inicialmente se estableció como objetivo principal del proyecto comparar varios algoritmos para encontrar el más óptimo, pero posteriormente se decidió comparar para diferentes algoritmo su desempeño en función de si la selección de datos fue manual o miente un algoritmo, además de desarrollar uno de los modelos para maximizar su desempeño.
    \item
    Escoger los entornos de desarrollo, de gestión y de comunicación.
    Como entorno de desarrollo vamos a utilizar Visual Sudio Code por la versatilidad que nos brinda, podremos desarrollar tanto la aplicación como el modelo en esta herramienta.
\end{itemize}
Este sprint se previó en una semana de duración pero realmente llevó a cabo en dos, para realizar una buena planificación y saber a dónde nos lleva el desarrollo del proyecto.

\subsection{Sprint 1: Comprensión del problema y buscar fuentes de datos (1 semana)} 
Este sprint se dividió en varias tareas:
\begin{itemize}
    \item
    Comprender el dominio de la F1.
    El objetivo principal de esta sección es obtener una comprensión profunda del mundo de las carreras de Fórmula 1 y familiarizarse con sus principales características. En este punto examinaré a fondo el dominio de este deporte de élite utilizando mi conocimiento y entusiasmo como aficionado a la Fórmula 1.
    A través de una investigación profunda de fuentes confiables y un análisis de datos históricos, además de mi experiencia personal como entusiasta de la Fórmula 1.
    \item
    Familiarizarse con las fuentes de datos disponibles y escoger una.
    Para este apartado cuento con mi experiencia como aficionado, ya que conozco muchas de las páginas que hablan de fórmula 1 y varias de ellas que albergan datos históricos. Inicialmente se han escogido para su estudio las siguientes webs: \href{https://www.statsf1.com/}{\textit{statsf1}}, \href{https://www.racing-reference.info/}{\textit{racing-reference}}, \href{https://www.f1-fansite.com/}{\textit{f1-fansite}},  \href{https://www.motorsportstats.com/}{\textit{motorsportstats}}, \href{https://www.formula1.com/}{\textit{Fórmula 1}} y \href{https://ergast.com/mrd/}{\textit{Ergast Developer API}}. 
    Una vez identificadas, se realizó un estudio de las fuentes de datos. Esto implicó examinar la precisión de los datos, la metodología aplicada para su recopilación, su consistencia y cualquier posible sesgo o limitación en el conjuntos total de datos. También, hay que conocer los formatos y estructuras de los datos disponibles. Algunas webs no permiten descargar los datos y hay que usar técnicas de webscrapping y es por eso que me he decantado por la API de Ergast, porque permite descargar los datos en formato json para un tratado de estos más fácil.
\end{itemize}

Este sprint se llevó a cabo en el tiempo esperado.

\subsection{Sprint 2: Descarga de datos (3 semanas)}
Para llevar a cabo el sprint se realizaron las siguientes tareas
\begin{itemize}
    \item 
    Creación de medios para la obtención de los datos.
    Una vez escogida la fuente se deben crear los medios necesarios para descargar esos datos. El objetivo de esta tarea fue crear clientes python para obtener los datos de la API de Ergast.
    Además, se obtuvieron datos de la página web de la página ofical de la Fórmula 1, que al final se descartaron por falta de información, lo que retrasó el proyecto algunos días. 
    \item
    Recopilación de los datos necesarios.
    En este momento se descargaron todos los datos suministrados por la API y además se decidió obtener algunos datos más para el entrenamiento, como puede ser la altitud del circuito con respecto al nivel del mar, el clima de las carreras, etc.
\end{itemize}

Este sprint se planificó para una semana pero realmente se tardaron 3 semanas debido a mi situación personal, estoy trabajando a tiempo completo y no dispuse de mucho tiempo para este proyecto, además del tiempo perdido tras la recopilación de datos en la página de la Fórmula 1.

\subsection{Sprint 3: Preparación de los datos (1 semanas)}
Para llevar a cabo la preparación de los datos se realizaron los siguientes pasos
\begin{itemize}
    \item 
    Preprocesar los datos para poder ser utilizados en la construcción del modelo.
    El objetivo de este paso fue preparar los datos descargados, entre esta preparación entra mapear los cambios de nombre de las escuderías para que no se tome a una escudería como dos diferentes, cambiar la nacionalidad de los pilotos por su país de procedencia y que coincida con el país de procedencia del equipo y dónde se encuentra el circuito entre otras cosas.
    \item
    Creación de nuevos datos a partir de los ya presentes,.
    Este paso trató de crear nuevas características para los entrenamientos o testeo de los modelos. Se crearón las siguientes nuevas características: fiabilidad de los coches, consistencia del piloto, edad del piloto, ganador de la carrera, ganador de la \textit{pole}, y se generó la clasifición del mundial de pilotos y equipos de cada año.
    \item 
    Limpieza de los datos.
    El objetivo de esta tarea fue limpiar los datos existentes, esto requiere eliminar o transformar los datos nulos. En este punto se decidió eliminar los datos nulos para que no afecten al entrenamiento.
    
\end{itemize}
Este sprint se llevo a cabo en el tiempo esperado, una semana.


\subsection{Sprint 4: Selección de características y codificación de los datos (1 semanas)}
Este sprint se dividió en dos tareas:
\begin{itemize}
\item 
    Codificación de los datos.
    Inicialmente se buscó codificar los datos, para ello se utilizaron varios algoritmos de codificación como One-Hot o codificación ordinal dependiendo del tipo de cada uno. Los datos que ya eran numéricos no se modificaron. Además, las fechas se codificaron en segundos en función de la fecha mínima de dicha características.
    \item 
    Selección de características.
    Esta tarea trató de generar dos selecciones, una manual y otra con un algoritmo de selección. Para ello se buscaron algoritmos y se decantó por el algoritmo Recursive Feature Elimination (RFE). Cómo tenemos tres variables objetivo se crearon cuatro conjuntos de datos, uno manual, y tres más diferentes según la variable objetivo.
\end{itemize}
La duración del sprint se estimó en una semana y se cumplió lo esperado.

\subsection{Sprint 5: Selección de los modelos, entrenamiento y evaluación de estos (2 semanas)}
Para llevar a cabo este sprint se realizaron dos tareas
\begin{itemize}
    \item Selección de modelos
    El objetivo fue elegir qué modelos se usaran para el entrenamiento. Se escogieron varios de ellos para que la comparación posterior sea más veraz.
    \item Entrenamiento y evaluación
    Esta tarea sirvió para entrenar los modelos iterativamente con cada uno de los datos escogidos y para cada una de las variables objetivo. Además se evalúa cada uno tras su entrenamiento guardando los datos en un diccionario para su posterior visionado.
\end{itemize}

\subsection{Sprint 6: Ajuste del modelo (3 semanas)}
Sprint dividido en 2 tareas:
\begin{itemize}
    \item 
    Seleccionar y recopilar los hiperparámetros de los modelos.
    La función de esta tarea fue escoger que modelos se pretendían ajustar y recopilar sus hiperparámetros más importantes. Se decidió ajustar el modelo Random Forest Regressor y el Decision Tree Regressor de la predicción del resultado de posivión en carrera.
    \item 
    Escoger los valores e iterar los entrenamientos.
    El objetivo fue escoger uno o varios valores para cada parñámetro e iterar su entrenamiento con el fin de lograr el mayor ejuste posible.
\end{itemize}


\section{Estudio de viabilidad}

Se va a separa este estudio en dos, el estudio de viabilidad legal y el estudio de viabilidad económica.

\subsection{Viabilidad económica}

Con el fin de una posible implementación del proyecto en un entorno real, desarrollaremos los costes y beneficios económicos estimados del proyecto en esta sección.

\subsubsection{Costes del proyecto}

Podemos dividir los costos del proyecto en los siguientes apartados: costos de personal, costos de hardware, costos de software, costos generales y beneficios potenciales.

\textbf{Costes personales}

Un ingeniero informático a tiempo completo trabajó en este proyecto durante dos meses. Se considera como el salario mínimo de un ingeniero \cite{salariogob}:

\begin{longtable}[]{@{}lr@{}}
\toprule
\begin{minipage}[b]{0.38\columnwidth}\raggedright\strut
\textbf{Concepto}\strut
\end{minipage} & \begin{minipage}[b]{0.20\columnwidth}\raggedright\strut
\textbf{Coste}\strut
\end{minipage}\tabularnewline
\midrule
\endhead
\begin{minipage}[t]{0.38\columnwidth}\raggedright\strut
Salario mensual neto\strut
\end{minipage} & \begin{minipage}[t]{0.20\columnwidth}\raggedright\strut
1.215,90\euro{}\strut
\end{minipage}\tabularnewline
\begin{minipage}[t]{0.38\columnwidth}\raggedright\strut
Retención IRPF (24\%)\strut
\end{minipage} & \begin{minipage}[t]{0.20\columnwidth}\raggedright\strut
633,01\euro{}\strut
\end{minipage}\tabularnewline
\begin{minipage}[t]{0.38\columnwidth}\raggedright\strut
Seguridad Social (29,9\%)\strut
\end{minipage} & \begin{minipage}[t]{0.20\columnwidth}\raggedright\strut
788,62\euro{}\strut
\end{minipage}\tabularnewline
\begin{minipage}[t]{0.38\columnwidth}\raggedright\strut
Salario mensual bruto\strut
\end{minipage} & \begin{minipage}[t]{0.20\columnwidth}\raggedright\strut
2.637,53\euro{}\strut
\end{minipage}\tabularnewline
\midrule
\begin{minipage}[t]{0.38\columnwidth}\raggedright\strut
\textbf{Total 2 meses}\strut
\end{minipage} & \begin{minipage}[t]{0.20\columnwidth}\raggedright\strut
5.275,06 \euro{}\strut
\end{minipage}\tabularnewline
\bottomrule
\caption{Costes de personal.}
\end{longtable}

Como estamos hablando del estudio de viabilidad económica del proyecto, para la cotización a la Seguridad Social se tienen en cuenta los tipos de cotización de la empresa. Estos valores han sido calculados utilizando los estándares establecidos por el gobierno, que podemos consultar en \cite{salariogob}. Estos son los diferentes formatos de cotización que se han utilizado:.

\begin{itemize}
\item
	Un 23,60\% por contingencias comunes.
\item
	Un 5,50\% por desempleo de tipo general.
\item
	Un 0,20\% para el Fondo de Garantía Salarial o FOGASA.
\item
	Y un 0,60\% para formación profesional.
\item
	Un total del 29,90\% del salario bruto total.
\end{itemize}

Para calcular la cotización del IRPF se ha consultado \cite{irpfgob}.\\

\textbf{Costes de hardware}

Hardware que será necesario para el desarrollo de este proyecto:

\begin{longtable}[]{@{}lr@{}}
\toprule
\begin{minipage}[b]{0.38\columnwidth}\raggedright\strut
\textbf{Hardware}\strut
\end{minipage} & \begin{minipage}[b]{0.20\columnwidth}\raggedright\strut
\textbf{Precio}\strut
\end{minipage}\tabularnewline
\midrule
\endhead
\begin{minipage}[t]{0.38\columnwidth}\raggedright\strut
Teclado y Ratón\strut
\end{minipage} & \begin{minipage}[t]{0.20\columnwidth}\raggedright\strut
20\euro{}\strut
\end{minipage}\tabularnewline
\begin{minipage}[t]{0.38\columnwidth}\raggedright\strut
Ordenador para el desarrollo\strut
\end{minipage} & \begin{minipage}[t]{0.20\columnwidth}\raggedright\strut
1000\euro{}\strut
\end{minipage}\tabularnewline
\midrule
\begin{minipage}[t]{0.38\columnwidth}\raggedright\strut
\textbf{Total}\strut
\end{minipage} & \begin{minipage}[t]{0.20\columnwidth}\raggedright\strut
1020\euro{}\strut
\end{minipage}\tabularnewline
\bottomrule
\caption{Costes de hardware.}
\end{longtable}

En el cálculo del precio se han seguido los siguientes criterios:

\begin{itemize}
    \item Precio del ordenador y periféricos usados en el desempeño de este proyecto.
\end{itemize}

\textbf{Costes de software}

Software utilizado en el desarrollo del proyecto:

\begin{itemize}
    \item Python 3.
    \item Visual Studio Code.
    \item Git.
\end{itemize}

Las licencias de software necesarias son de código abierto o gratuitas, y dado que se pueden ejecutar en Linux, un sistema operativo de código abiertos, no hay ningún costo por el software utilizado.

\textbf{Costes de software}

Costes totales del proyecto:

\begin{longtable}[]{@{}lr@{}}
\toprule
\begin{minipage}[b]{0.38\columnwidth}\raggedright\strut
\textbf{Coste}\strut
\end{minipage} & \begin{minipage}[b]{0.20\columnwidth}\raggedright\strut
\textbf{Valor}\strut
\end{minipage}\tabularnewline
\midrule
\endhead
\begin{minipage}[t]{0.38\columnwidth}\raggedright\strut
Coste personal\strut
\end{minipage} & \begin{minipage}[t]{0.20\columnwidth}\raggedright\strut
5.275,06\euro{}\strut
\end{minipage}\tabularnewline
\begin{minipage}[t]{0.38\columnwidth}\raggedright\strut
Coste de hardware\strut
\end{minipage} & \begin{minipage}[t]{0.20\columnwidth}\raggedright\strut
1.020\euro{}\strut
\end{minipage}\tabularnewline
\begin{minipage}[t]{0.38\columnwidth}\raggedright\strut
Coste software\strut
\end{minipage} & \begin{minipage}[t]{0.20\columnwidth}\raggedright\strut
0\euro{}\strut
\end{minipage}\tabularnewline
\midrule
\begin{minipage}[t]{0.38\columnwidth}\raggedright\strut
\textbf{Total}\strut
\end{minipage} & \begin{minipage}[t]{0.20\columnwidth}\raggedright\strut
6.295,06\euro{}\strut
\end{minipage}\tabularnewline
\bottomrule
\caption{Costes de totales del proyecto.}
\end{longtable}

\textbf{Beneficios}

En cuanto a los beneficios que podría aportar el desarrollo del proyecto, podríamos vender la solución a un equipo del campeonato para que puedan usarlo, o a una casa de apuestas para establecer los multiplicadores de acierto.




\apendice{Especificación de Requisitos}

\section{Introducción}

Una descripción detallada del comportamiento del sistema a desarrollar es lo que se pretende lograr con la Especificación de requisitos de software (ERS). Además de delinear los requisitos del sistema, también enumera las necesidades del producto (tanto para el usuario como para el cliente).
Cubre todos los casos posibles que un usuario final puede realizar con el software.

Este documento sirve de medio de comunicación entre todas las partes (desarrolladores, clientes y usuarios finales).

Según el estándar IEEE-STD-830-1998 \cite{ieee830} una ERS debe presentar las siguientes características:

\begin{itemize}
    \item \textbf{Correcto}: es correcto si, y sólo si, el software satisface todos los requisitos especificados.
    \item \textbf{Consistente}: debe ser coherente con los requerimientos de la ERS, así como con los documentos de distinto nivel.
    \item \textbf{Inequívoco}: es inequívoco si, y sólo si, solo hay una interpretación de cada requisito establecido.
    \item \textbf{Completo}: para ello debe incluir: 
        \begin{itemize}
            \item Requisitos relacionados a desarrollo, funcionalidad, restricciones de diseño, atributos e interfaces externas.
            \item Definiciones de respuestas de software a todos los datos de entrada y a todas las circunstancias posibles.
            \item Tener las etiquetas llenas y bien referencias. 
        \end{itemize}
    \item \textbf{Delinear que tiene estabilidad y/o importancia}: La importancia y/o estabilidad de cada requisito debe describirse en una ERS. Cada requisito es poseedor de un identificador único que indica su estabilidad o importancia.
    \item \textbf{Comprobable}: cada requisito declarado debe ser verificable. Si existe una forma concreta de verificar que el producto de software cumple con el requisito, entonces ese requisito es verificable.
    \item \textbf{Modificable}: si su estructura y estilo permiten que los cambios en los requisitos se realicen de manera rápida, completa y consistente mientras se conserva la estructura y estilo.
    \item \textbf{Identificable}: una ERS es identificable si la fuente de los requisitos es obvia y facilita la referencia en un posible desarrollo posterior o la documentación del mismo.
\end{itemize}

\section{Objetivos generales}

Se pretende cumplir los objetivos siguientes:

\begin{itemize}
\tightlist
\item
  Desarrollar un modelo para la predicción de resultados en las carreras de Fórmula 1.
\item
  Comparar una selección mediante algoritmos contra una selección manual.
\item
  Ajustar un modelo para mejorar sus predicciones.
\item
  Desarrollar una aplicación para  la interacción con el modelo final.
\end{itemize}


\section{Catálogo de requisitos}

En esta sección se presentan los requisitos del sistema tanto funcionales como no funcionales.

\subsection{Requisitos funcionales}

\begin{itemize}
    \item \textbf{RF-1 Gestión del modelo entrenado:} 
        El sistema debe de ser capaz de realizar predicciones.
        \begin{itemize}
            \item \textbf{RF-1.1 Predicciones de resultados:}
                El modelo debe poder predecir los resultados.
                \begin{itemize}
                    \item \textbf{RF-1.1.1 Predicciones de resultados de pilotos:}
                        El modelo debe poder predecir los resultados de los pilotos en una carrera determinada.
                    \item \textbf{RF-1.1.2 Predicciones del ganador:}
                        El modelo debe ser capaz de predecir si un piloto ganará o no una carrera.
                    \item \textbf{RF-1.1.3 Predicciones del poleman:}
                        El modelo debe ser capaz de predecir si un piloto hará o no la pole.
				\end{itemize}
        \end{itemize}
    \item \textbf{RF-2 Gestión de la Interfaz de usuario:}
        La interfaz de usuario debe ser capaz de realizar las funciones necesarias para interactuar con el modelo de forma simple y fácil para el usuario.
        \begin{itemize}
            \item \textbf{RF-2.1 Predicciones de resultados:}
                La aplicación debe poder predecir los resultados.
                \begin{itemize}
                    \item \textbf{RF-2.1.1 Predicciones de resultados de pilotos:}
                        La aplicación debe poder predecir los resultados de los pilotos en una carrera determinada.
                    \item \textbf{RF-2.1.2 Predicciones del ganador:}
                        La aplicación debe ser capaz de predecir si un piloto ganará o no una carrera.
                    \item \textbf{RF-2.1.3 Predicciones del poleman:}
                        La aplicación debe ser capaz de predecir si un piloto hará o no la pole.
				\end{itemize}
        \end{itemize}
\end{itemize}
    


\subsection{Requisitos no funcionales}

\begin{itemize}
    \item \textbf{RNF-1 Usabilidad:} la aplicación debe mantener una interfaz de usuario intuitiva y ser lo suficientemente similar a otras aplicaciones del mismo entorno para que no sea necesario un tutorial.
    \item \textbf{RNF-2 Rendimiento:} independientemente del dispositivo utilizado, tanto el modelo como la aplicación deben permanecer fluidos, sin fallos ni paradas en el sistema.
    \item \textbf{RNF-3 Disponibilidad:} el modelo debe ser accesible en todo momento y además, la aplicación debe estar disponible para su uso independientemente del momento.
    \item \textbf{RNF-4 Estabilidad:} el sistema debe ser estable, manteniendo un bajo nivel de errores y poder ocultarlos al usuario.
    \item \textbf{RNF-5 Robustez:} los eventos inesperados no deben causar un bloqueo o falla del sistema.
    \item \textbf{RNF-6 Fiabilidad (validez e integridad):} todas las acciones realizadas dentro del sistema deben ser válidas, preservando la integridad de los datos.
    \item \textbf{RNF-7 Mantenibilidad:} el sistema debe admitir la mejora del rendimiento (escalabilidad), la corrección de fallas y la adaptación de una manera simple.
    \item \textbf{RNF-8 Portabilidad:} el modelo debe poder ser ejecutado en varios entornos o plataformas, lo que le permita funcionar en diferentes sistemas operativos o configuraciones de hardware.
\end{itemize}



\section{Especificación de requisitos}

En este apartado se muestran los casos de uso derivados de los requisitos funcionales del proyecto.

\subsection{Caso de uso principal}

\imagen{caso-general}{Caso de uso principal de la especificación de requisitos.}

\subsection{Actores}

Un único actor, usuario final que actúa con el sistema.

\subsection{Casos de uso}


\apendice{Especificación de diseño}

\section{Introducción}

Este anexo detalla los datos que maneja el sistema, la arquitectura y su diseño procedimental para explicar cómo se han cumplido los requisitos y objetivos del proyecto.

\section{Diseño de datos}

En esta sección, se describe cómo se han recopilado, procesado y estructurado los datos relacionados con las carreras de Fórmula 1.

\subsection{Recopilación}

La mayor parte de los datos están recopilados de una API llamada \href{https://ergast.com/mrd/}{\textit{Ergast Developer API}}, que es un servicio web experimental que proporciona un registro histórico de datos de carreras de coches para fines no comerciales. Se han obtenido de la API los siguientes datos: pilotos, circuitos, equipos y resultados de las carreras. Para ello se han ejecutado y obtenido los siguientes archivos:
\begin{enumerate}
    \item Pilotos: debemos ejecutar el archivo \texttt{1\_GetDrivers} para obtener los datos de los pilotos. Una vez ejecutados se guardan el la capeta \textit{code/data} con el nombre \texttt{Drivers.csv}.
    \item Circuitos: debemos ejecutar el archivo \texttt{2\_GetCircuits} para obtener los datos de los circuitos. Una vez ejecutados se guardan el la capeta \textit{code/data} con el nombre \texttt{Circuits.csv}.
    \item Equipos: debemos ejecutar el archivo \texttt{3\_GetConstructors} para obtener los datos de los equipos. Una vez ejecutados se guardan el la capeta \textit{code/data} con el nombre \texttt{Constructors.csv}.
    \item Resultados de las carreras: debemos ejecutar el archivo \texttt{4\_GetResults} para obtener los datos de las carreras. Una vez ejecutados se guardan el la capeta \textit{code/data} con el nombre \texttt{Results.csv}.
    \item Clasificaciones: también se pueden obtener datos de las clasificaciones de esta API, pero como no se han usado, no se han guardado, En caso de necesitarlo se debe ejecutar el archivo \texttt{6\_GetQualyDataErgast}. Una vez ejecutados se guardarían el la capeta \textit{code/data} con el nombre \texttt{qualifying.csv}.
\end{enumerate}

Además, se han obtenido datos del clima de las carreras de la página de la Wikipedia a través de webscrapping, para ello ejecutaremos el archivo \texttt{5\_GetWeatherData} para obtener los datos meteorológicos. Una vez ejecutado se añade en los datos de los resultados y se guardan el la capeta \textit{code/data} con el nombre \texttt{final\_results.csv}.

\subsection{Procesado}

Para procesar los datos se ha seguido la siguiente estrategia:

\begin{enumerate}
    \item Combinación de los datos: se han combinado los datos de los archivos: \texttt{Drivers.csv}, \texttt{Circuits.csv}, \texttt{Constructors.csv} y \texttt{Results.csv}. 
    \item Preparación de los datos: a partir de los datos combinados se han creado datos de finalización de carrera, consistencia de pilotos, fiabilidad de los coches, país de procedencia de equipos y pilotos, edad de los pilotos durante las carreras, ganadores y polemans, además de corregir algunos datos como pueden ser los cambios de nombre del mismo equipo en determinados años. Todo ello ejecutando los archivos \texttt{7\_A\_PrepareAndMergeData} y \texttt{7\_B\_DriverConstructorData}, y guardado en el archivo \texttt{merged\_data.csv} en la carpeta \textit{code/data}.
    \item Limpieza de datos: a partir de los datos obtenidos en \texttt{merged\_data.csv} se han eliminado las características que contenían datos nulos a través del archivo \texttt{8\_A\_CleanData}, y guardado en el archivo \texttt{all\_data.csv} en la carpeta \textit{code/data}.
    \item Obtención de las clasificaciones de pilotos y equipos: a partir de los datos obtenidos en \texttt{all\_data.csv} se han generado los archivos \texttt{driver\_standings.csv} y \texttt{constructor\_standings.csv} guardados en la carpeta \textit{code/data}.
    \item Codificación y selección de los datos: se han codificado los datos para entrenarlos en los modelos a través de los archivos: \texttt{9\_A\_ManualSelection}, \texttt{9\_B\_CodificationManualData}, \texttt{10\_A\_CodificationData}, \texttt{10\_B\_AutoSelection}, y resultando en los siguientes archivos de datos guardados en la carpeta \textit{code/data}:
        \begin{itemize}
            \item \texttt{coded\_auto\_selection\_data.csv}: datos codificados para la selección de carcaterísticas con el algoritmo RFE.
            \item \texttt{coded\_auto\_selection\_data\_race\_winner.csv}: datos con las características ya seleccionadas para la variable objetivo ganador de la carrera.
            \item \texttt{coded\_auto\_selection\_data\_race\_final\_position.csv}: datos con las características ya seleccionadas para la variable objetivo posiciones finales de carrera.
            \item \texttt{coded\_auto\_selection\_data\_qualy\_pole.csv}: datos con las características ya seleccionadas para la variable objetivo ganador de la pole.
            \item \texttt{coded\_manual\_data.csv}: datos seleccionados manualmente.
            \item \texttt{codificadores.pkl}: se han guardado lo codificadores utilizados para cada característica en la carpeta \textit{code/coders}.
        \end{itemize}
    \item Emtrenamiento de los modelos: los modelos han sido entrenados en el archivo \texttt{Training}, el cual ha generado los archivos con las características más relevantes para el ajuste posterior en la carpeta \textit{code/best\_split\_data}. Además, se han guardado los modelos entrenados en la carpeta \textit{code/train\_models}.
    \item Ajuste de los modelos: encontramos los modelos con ajuste en el archivo \texttt{TrainingParams}, que ha guardado el modelo con mejor ajuste en la carpeta \textit{code/train\_models}.
\end{enumerate}

\subsection{Aplicación}

Cómo la aplicación es de interacción con el modelo entrenado, no guarda una persistencia de los datos como tal. Contamos con 3 archivos de datos en los que se guardan circuitos constructores y pilotos de la temporada 2023 de Fórmula 1 para hacer las predicciones:
\begin{itemize}
    \item Pilotos: \texttt{app\_drivers}.
    \item Constructores: \texttt{app\_constructors}.
    \item Circuitos: \texttt{app\_circuits}.
\end{itemize}

\section{Diseño procedimental}

Podemos entender las operaciones internas del sistema gracias a su diseño procedimental.

La interacción entre los diversos componentes del sistema se representará mediante diagramas de secuencia teniendo en cuenta los requisitos del sistema.

\imagen{seqpredpos}{Diagramas de secuencia relativo al RF-2.1.1.}
\imagen{seqpredwinner}{Diagramas de secuencia relativo al RF-2.1.2.}
\imagen{seqpredpole}{Diagramas de secuencia relativo al RF-2.1.3.}

\section{Diseño arquitectónico}

\subsection{Arquitectura del sistema}

La arquitectura modelo-vista-controlador (MVC) es el patrón de diseño que se ha seguido en la aplicación.

\subsubsection{Patrón modelo-vista-controlador (MVC)}

El patrón MVC (figura \ref{fig:mvc}) es un patrón de diseño arquitectónico que divide el sistema en tres componentes: la interfaz de usuario, la lógica de la aplicación y los datos. La capa de datos está representada por el modelo, la interfaz de usuario está representada por la vista y la lógica de la aplicación está representada por el controlador.

\imagen{mvc}{Explicación del patrón MVC.}

En este proyecto se ha utilizado de la siguiente manera:
\begin{itemize}
    \item Una primera carpeta de datos (\textit{data}) que contiene los datos necesarios para hacer las predicciones.
    \item Una carpeta llamada \textit{logic} que contiene la lógica de la aplicación.
    \item Una última carpeta que implementa la capa de presentación (\textit{presentation}).
\end{itemize}

\subsection{Estructura global}

Por último podemos ver el diagrama que componen la aplicación de interacción con el modelo en la figura \ref{fig:arqmvc}.

\imagen{arqmvc}{Diagrama de clases de la aplicación.}

\apendice{Documentación técnica de programación}

\section{Introducción}

\section{Estructura de directorios}

\section{Manual del programador}

\section{Compilación, instalación y ejecución del proyecto}

\section{Pruebas del sistema}

\apendice{Documentación de usuario}

\section{Introducción}

\section{Requisitos de usuarios}

\section{Instalación}

\section{Manual del usuario}





\bibliographystyle{plain}
\bibliography{bibliografiaAnexos}

\end{document}
